\section{Statistical Precision as Stiffness and Friction}\label{sec:mass}

The natural-gradient flow of Eq.~\eqref{eq:gauge_natural_gradient} descends the consensus free-energy functional $\mathcal{F}$ of Section~\ref{sec:free_energy}, which is therefore the potential of the agent dynamics. This section computes the second variation of $\mathcal{F}$ in the Gaussian belief sector and derives the two roles of statistical precision in that first-order dynamics (status D, Table~\ref{tab:claim_status}). The Hessian is a stiffness tensor whose blocks are built from precisions: the prior term contributes the prior precision $\bar{\Lambda}_{p_i}$, and the interaction terms contribute transported neighbor precisions. The Fisher metric, which in the mean sector is the belief precision, is the friction tensor of the flow, and no term of the dynamics carries an inertia. The only velocity-quadratic form in the dynamics is the Rayleigh dissipation function $\mathcal{R}$, which is half the free-energy dissipation rate and fixes the rate of the information clock of Section~\ref{sec:time_info} through $(d\vartheta/dt)^2 = 2\eta\mathcal{R}$ (Section~\ref{sec:stiffness_friction}). An inertial reading of precision needs a separate kinetic-metric postulate with its own second-order equation of motion, stated in Section~\ref{sec:velocity_quadratic}.

\paragraph{Stiffness, friction, and inertia.} For a damped oscillator $m\ddot{x} + \gamma\dot{x} + kx = 0$ the three coefficients are independent data: the Hessian of the potential gives the stiffness $k$, the dissipation gives the friction $\gamma$, and the kinetic energy gives the mass $m$. The Hessian of $\mathcal{F}$ derived below supplies $k$, the Fisher metric of the natural-gradient flow supplies $\gamma$, and neither $\mathcal{F}$ nor the flow supplies $m$. We hold the prior parameters $(\bar{\mu}_i, \bar{\Sigma}_i)$ fixed, the quasi-static approximation in which they evolve slowly relative to the beliefs $(\mu_i, \Sigma_i)$.

\subsection{Setup and Notation}

Each agent $i$ maintains a belief distribution $q_i = \mathcal{N}(\mu_i, \Sigma_i)$ anchored to a prior $p_i = \mathcal{N}(\bar{\mu}_i, \bar{\Sigma}_i)$ and receives observations $o_i$ through the identity readout $p(o_i \mid k_i, m_i) = \mathcal{N}(o_i; k_i, \Sigma_{o_i})$, whose likelihood does not depend on the model parameter $m_i$ and is written $p(o_i \mid k_i)$ below. We define precision matrices
\begin{align}
\Lambda_{q_i} &= \Sigma_i^{-1} & &\text{(belief precision)} \\
\bar{\Lambda}_{p_i} &= \bar{\Sigma}_i^{-1} & &\text{(prior precision)} \\
\Lambda_{o_i} &= \Sigma_{o_i}^{-1} & &\text{(observation precision)},
\end{align}
with transported quantities $\tilde{\mu}_k = \Omega_{ik}\mu_k$ and $\tilde{\Sigma}_k = \Omega_{ik}\Sigma_k\Omega_{ik}^\top$, where $\Omega_{ik} \in \mathrm{GL}(K_q)$ is the transport of the interaction term: in Regime I the frame quotient $e^{\phi_i}e^{-\phi_k}$, and in Regime II an independent link on the interaction graph with orientation reversal $\Omega_{ki} = \Omega_{ik}^{-1}$. The block formulas below use only the transports and hold in both regimes. The transported precision is the inverse of the transported covariance,
\begin{equation}
\tilde{\Lambda}_{q_k} := (\Omega_{ik}\Sigma_k\Omega_{ik}^\top)^{-1} = \Omega_{ik}^{-\top}\Lambda_{q_k}\Omega_{ik}^{-1},
\label{eq:precision_transport}
\end{equation}
which differs from the covariance transport law because precision is the inverse of covariance. For orthogonal transport $\Omega \in \mathrm{O}(K_q)$ the two laws coincide ($\Omega^{-\top} = \Omega$), and Eq.~\eqref{eq:precision_transport} reduces to $\Omega\Lambda\Omega^\top$; for general $\mathrm{GL}(K_q)$ the dual transport is essential, and we use Eq.~\eqref{eq:precision_transport} throughout. The boxed identity $\Omega_{ji}^\top \tilde{\Lambda}_{q_i}^{(j)} \Omega_{ji} = \Lambda_{q_i}$ used below holds in $\mathrm{GL}(K_q)$ because $\Omega^\top \Omega^{-\top} = I$ algebraically; the secondary simplifications $\tilde{\Lambda}_k\Omega_{ik} = \Omega_{ik}\Lambda_k$ and $\Omega_{ki}^\top\tilde\Lambda_i = \Lambda_i\Omega_{ki}^\top$ hold only under $\mathrm{O}(K_q)$; the general-linear identity is given below.

\subsection{The Extended Free Energy Functional}

The extended working functional at one base point, the fixed-$\beta^*$ block-coordinate form of Eq.~\eqref{eq:free_energy_functional_final} with explicit sensory evidence (the attention-entropy block is omitted as constant in the beliefs; this is not the entropy-suppressed surrogate of Section~\ref{sec:complete_free_energy}), is
\begin{equation}
\mathcal{F}[\{q_i\}] = \sum_i D_{\mathrm{KL}}(q_i \| p_i) + \sum_{i,k} \beta_{ik} D_{\mathrm{KL}}(q_i \| \Omega_{ik}[q_k]) - \sum_i \mathbb{E}_{q_i}[\log p(o_i \mid k_i)].
\label{eq:extended_free_energy}
\end{equation}
The three terms represent prior anchoring (deviation from the prior $p_i$, which in the hierarchy of Section~\ref{sec:cross_scale_shadows} is the one-episode-lagged predictive prior supplied by the parent level), interaction (alignment with other agents through the transports), and sensory evidence. Under the product recognition law $q_i \otimes s_i$ and the identity readout, the jointly typed observation term $-\mathbb{E}_{q_i \otimes s_i}[\log p(o_i \mid k_i, m_i)]$ of Eq.~\eqref{eq:free_energy_functional_final} reduces to the $q_i$-expectation shown, and at fixed $s_i$ the model-fiber terms are constant in $(\mu_i, \Sigma_i)$ and drop out of every variation below. Under the frozen-$\beta^*$ convention adopted for this section (the differentiation-convention paragraph below), the omitted attention-entropy term $\tau\beta_{ik}\log(\beta_{ik}/\pi_{ik})$ of the canonical functional is an additive constant in $(\mu, \Sigma)$, so the first and second variations taken under that convention coincide with the frozen-$\beta$ variations of Eq.~\eqref{eq:free_energy_functional_final}.

\subsection{Component Free Energies for Gaussians}

For $q = \mathcal{N}(\mu_q, \Sigma_q)$ and $p = \mathcal{N}(\mu_p, \Sigma_p)$, the relative entropy is given by the closed-form expression of Eq.~\eqref{eq:gaussian_kl}, applied with the substitution $(\mu_1, \Sigma_1, \mu_2, \Sigma_2) \to (\mu_q, \Sigma_q, \mu_p, \Sigma_p)$.

For the Gaussian likelihood $p(o_i \mid k_i) = \mathcal{N}(o_i; k_i, \Sigma_{o_i})$, the expected log-likelihood evaluates to
\begin{equation}
-\mathbb{E}_{q_i}[\log p(o_i \mid k_i)] = \frac{1}{2}(o_i - \mu_i)^\top \Lambda_{o_i}(o_i - \mu_i) + \frac{1}{2}\operatorname{tr}(\Lambda_{o_i}\Sigma_i) + \mathrm{const}.
\end{equation}

\subsection{First Variations (Gradient)}

\paragraph{Prior Term}
\begin{align}
\frac{\partial D_{\mathrm{KL}}(q_i \| p_i)}{\partial \mu_i} &= \bar{\Lambda}_{p_i}(\mu_i - \bar{\mu}_i), \\
\frac{\partial D_{\mathrm{KL}}(q_i \| p_i)}{\partial \Sigma_i} &= \frac{1}{2}(\bar{\Lambda}_{p_i} - \Lambda_{q_i}).
\end{align}

\paragraph{Interaction Term}

With respect to receiver $i$:
\begin{align}
\frac{\partial D_{\mathrm{KL}}(q_i \| \tilde{q}_k)}{\partial \mu_i} &= \tilde{\Lambda}_{q_k}(\mu_i - \tilde{\mu}_k), \\
\frac{\partial D_{\mathrm{KL}}(q_i \| \tilde{q}_k)}{\partial \Sigma_i} &= \frac{1}{2}(\tilde{\Lambda}_{q_k} - \Lambda_{q_i}).
\end{align}

With respect to sender $k$ (using the general-linear precision transport of Eq.~\eqref{eq:precision_transport}):
\begin{align}
\frac{\partial D_{\mathrm{KL}}(q_i \| \tilde{q}_k)}{\partial \mu_k} &= \Lambda_{q_k}\Omega_{ik}^{-1}(\tilde{\mu}_k - \mu_i), \\
\frac{\partial D_{\mathrm{KL}}(q_i \| \tilde{q}_k)}{\partial \Sigma_k} &= \frac{1}{2}\Omega_{ik}^\top\left[\tilde{\Lambda}_{q_k} - \tilde{\Lambda}_{q_k}\big(\Sigma_i + \tilde{d}_{ik}\tilde{d}_{ik}^\top\big)\tilde{\Lambda}_{q_k}\right]\Omega_{ik},
\end{align}
where $\tilde{d}_{ik} := \tilde{\mu}_k - \mu_i = \Omega_{ik}\mu_k - \mu_i$ is the receiver-frame mean residual. The quadratic-form contribution $-\tilde{\Lambda}_{q_k}\tilde{d}_{ik}\tilde{d}_{ik}^\top\tilde{\Lambda}_{q_k}$ originates in differentiating the Mahalanobis term $\tilde{d}_{ik}^\top \tilde{\Lambda}_{q_k}\tilde{d}_{ik}$ with respect to $\tilde{\Sigma}_k$ via $\partial \tilde{\Lambda}_{q_k}/\partial \tilde{\Sigma}_k = -\tilde{\Lambda}_{q_k}\otimes\tilde{\Lambda}_{q_k}$, then chain-ruled through $\tilde{\Sigma}_k = \Omega_{ik}\Sigma_k\Omega_{ik}^\top$. The outer $\Omega_{ik}^\top(\cdot)\Omega_{ik}$ structure follows from the same chain rule applied to the leading $\tilde{\Lambda}_{q_k}$ and $\tilde{\Lambda}_{q_k}\Sigma_i\tilde{\Lambda}_{q_k}$ contributions; for the mean gradient, the same chain rule yields the explicit inverse $\Omega_{ik}^{-1}$, which for orthogonal $\Omega$ would equal $\Omega_{ik}^\top$. The $\tilde{d}_{ik}\tilde{d}_{ik}^\top$ term vanishes at consensus ($\tilde{d}_{ik} = 0$), so the boxed stiffness blocks of Eqs.~\eqref{eq:mass_mu_diagonal}--\eqref{eq:mass_sigma_offdiagonal}, evaluated at consensus, are unchanged. Off consensus the term contributes a rank-one stiffness aligned with $\tilde{d}_{ik}$, so any second-variation analysis away from the consensus point must retain it.

\paragraph{Sensory Term}
\begin{align}
\frac{\partial}{\partial \mu_i}\left[-\mathbb{E}_{q_i}[\log p(o_i \mid k_i)]\right] &= \Lambda_{o_i}(\mu_i - o_i), \\
\frac{\partial}{\partial \Sigma_i}\left[-\mathbb{E}_{q_i}[\log p(o_i \mid k_i)]\right] &= \frac{1}{2}\Lambda_{o_i}.
\end{align}

\paragraph{Total Gradient}
\begin{equation}
\frac{\partial \mathcal{F}}{\partial \mu_i} = \bar{\Lambda}_{p_i}(\mu_i - \bar{\mu}_i) + \sum_k \beta_{ik}\tilde{\Lambda}_{q_k}(\mu_i - \tilde{\mu}_k) + \sum_j \beta_{ji}\Lambda_{q_i}\Omega_{ji}^{-1}(\tilde{\mu}_i^{(j)} - \mu_j) + \Lambda_{o_i}(\mu_i - o_i),
\label{eq:total_gradient_mu}
\end{equation}
where $\tilde{\mu}_i^{(j)} = \Omega_{ji}\mu_i$ is agent $i$'s mean transported into agent $j$'s frame.

\subsection{Second Variations: The Stiffness Tensor}

The frozen-$\beta$ Hessian of the free energy, $M = \partial^2 \mathcal{F}/\partial\xi\partial\xi$ with the attention weights held at $\beta^*$ (see the differentiation-convention paragraph below), is a symmetric bilinear form on belief configuration space, the stiffness tensor of the potential $\mathcal{F}$. It is positive definite in the covariance sector at the reciprocal-consensus point established below, and it is not in general a metric: the reduced-functional correction and the state-dependent prior weight discussed below both add negative-semidefinite terms away from consensus. The Hessian is taken in the flat affine chart $(\mu_i, \Sigma_i) \in \mathbb{R}^{K_q} \times \mathrm{Sym}(K_q)$. A coordinate Hessian is a tensor, independent of the chart, only at a critical point of $\mathcal{F}$, and consensus is a critical point only when the prior and sensory terms are also stationary there. Away from critical points the chart matters: for $f(\Sigma) = D_{\mathrm{KL}}(\mathcal{N}(0,\Sigma)\|\mathcal{N}(0,\Sigma_p)) + \tfrac12\operatorname{tr}(\Lambda_o\Sigma)$ the second variations along the same tangent vectors in the $\Sigma$ chart and in the precision chart $\Lambda = \Sigma^{-1}$ agree at the critical point but differ by $49$ percent at a noncritical point. A chart-free alternative is the covariant Hessian of the Fisher Levi-Civita connection, which differs from the coordinate Hessian by a Christoffel-symbol term proportional to the gradient and therefore coincides with it at critical points. The full state vector is $\xi = (\mu_1, \ldots, \mu_N, \Sigma_1, \ldots, \Sigma_N)$, yielding the block structure
\begin{equation}
M = \begin{pmatrix} M_{\mu\mu} & C^{\mu\Sigma} \\ (C^{\mu\Sigma})^\top & M_{\Sigma\Sigma} \end{pmatrix},
\label{eq:mass_block_structure}
\end{equation}
where each block is itself an $N \times N$ matrix of sub-blocks.

\paragraph{Notation hierarchy.}
We use three related but distinct objects throughout: (i) $M_{\mu\mu}$ (a matrix), the full Hessian of the mean sector defined in \eqref{eq:mass_block_structure}; (ii) $\bar{\Lambda}_{p_i} = \bar{\Sigma}_i^{-1}$ (a matrix), the prior precision of agent $i$, which is the whole of $[M_{\mu\mu}]_{ii}$ when inter-agent coupling and observation likelihoods are absent ($\beta_{ij}=0$, $\Lambda_{o_i}=0$); and (iii) the scalar stiffness $\bar{M}_i := \operatorname{tr}([M_{\mu\mu}]_{ii})/K_q$, used when a one-number summary is needed. When unqualified, ``the stiffness tensor'' refers to (i) and ``stiffness'' as a scalar refers to (iii).

\paragraph{Differentiation convention.}
The second-variation expressions below treat the attention weights $\beta_{ij}$ as fixed at their softmax-equilibrium values $\beta_{ij}^*$. The envelope theorem licenses this frozen-$\beta$ convention at first order only: for the reduced free energy with $\beta$ eliminated, the true Hessian adds the correction $-\tau^{-1}\mathrm{Cov}_{\beta^*}(\nabla E_{ij}, \nabla E_{ij})$ to the frozen-$\beta$ blocks (the reduced-Hessian identity~\eqref{eq:rg_reduced_hessian} of Appendix~\ref{sec:meta_agent_rg}), a negative-semidefinite term that vanishes exactly at consensus, where the per-edge gradients $\nabla E_{ij}$ coincide, but persists at generic equilibria in which prior, interaction, and sensory pulls balance with nonzero per-edge gradients, so the frozen-$\beta$ blocks overstate the curvature of the reduced functional away from consensus. Differentiating through the softmax rows produces these cross terms through $\partial \beta_{ij}/\partial \mu$ and $\partial \beta_{ij}/\partial \Sigma$. The correction is bounded: writing $D_i = \max_{j,k}\|\nabla E_{ij} - \nabla E_{ik}\|$ for the per-edge gradient diameter, the identity $\mathrm{Cov}_\beta(X) = \tfrac{1}{2}\sum_{j,k}\beta_j\beta_k(X_j - X_k)(X_j - X_k)^\top$ gives the spectral bound $\|\tau^{-1}\mathrm{Cov}_{\beta^*}(\nabla E_{ij})\| \le D_i^2/(2\tau)$. At nondegenerate energy gaps the covariance decays exponentially in the gap over $\tau$, so the correction vanishes in the sharp-attention limit away from ties and is large precisely near attention switching points, where per-edge energies nearly coincide at small $\tau$; this is the rapid-reorganization regime in which the frozen-$\beta$ blocks are unreliable. The boxed forms below are written for general $\beta_{ij}^* \geq 0$ under this convention and reduce to the isolated-agent limit by setting $\beta_{ij}^* = 0$ in the sums. The first variations of the preceding subsection need no such convention, since by the envelope theorem the frozen-$\beta$ gradient equals the gradient of the reduced functional at $\beta = \beta^*$; the convention is stated here because it first matters at second order.

\paragraph{Mean Sector}

\paragraph{Diagonal blocks $(i = k)$:} From the prior we obtain $\partial^2 D_{\mathrm{KL}}(q_i \| p_i)/\partial \mu_i \partial \mu_i^\top = \bar{\Lambda}_{p_i}$. From the interaction term as receiver, $\partial^2 D_{\mathrm{KL}}(q_i \| \tilde{q}_k)/\partial \mu_i \partial \mu_i^\top = \tilde{\Lambda}_{q_k}$. From the interaction term as sender to agent $j$, $\partial^2 D_{\mathrm{KL}}(q_j \| \tilde{q}_i)/\partial \mu_i \partial \mu_i^\top = \Omega_{ji}^\top \tilde{\Lambda}_{q_i}^{(j)} \Omega_{ji} = \Lambda_{q_i}$. From sensory evidence, $\partial^2[-\mathbb{E}_{q_i}[\log p(o_i \mid k_i)]]/\partial \mu_i \partial \mu_i^\top = \Lambda_{o_i}$. The total diagonal stiffness block is
\begin{equation}
\boxed{[M_{\mu\mu}]_{ii} = \bar{\Lambda}_{p_i} + \sum_k \beta_{ik} \tilde{\Lambda}_{q_k} + \sum_j \beta_{ji} \Lambda_{q_i} + \Lambda_{o_i}.}
\label{eq:mass_mu_diagonal}
\end{equation}

\paragraph{Off-diagonal blocks $(i \neq k)$:} From $D_{\mathrm{KL}}(q_i \| \tilde{q}_k)$, using Eq.~\eqref{eq:precision_transport},
\begin{equation}
\frac{\partial^2 D_{\mathrm{KL}}(q_i \| \tilde{q}_k)}{\partial \mu_i \partial \mu_k^\top} = -\tilde{\Lambda}_{q_k}\Omega_{ik} = -\Omega_{ik}^{-\top}\Lambda_{q_k},
\end{equation}
where the second equality follows from $\tilde{\Lambda}_{q_k}\Omega_{ik} = \Omega_{ik}^{-\top}\Lambda_{q_k}\Omega_{ik}^{-1}\Omega_{ik} = \Omega_{ik}^{-\top}\Lambda_{q_k}$. From $D_{\mathrm{KL}}(q_k \| \tilde{q}_i)$ (if $k$ also attends to $i$),
\begin{equation}
\frac{\partial^2 D_{\mathrm{KL}}(q_k \| \tilde{q}_i)}{\partial \mu_i \partial \mu_k^\top} = -\Lambda_{q_i}\Omega_{ki}^{-1}.
\end{equation}
The sensory term does not couple different agents. Therefore
\begin{equation}
\boxed{[M_{\mu\mu}]_{ik} = -\beta_{ik}\Omega_{ik}^{-\top}\Lambda_{q_k} - \beta_{ki}\Lambda_{q_i}\Omega_{ki}^{-1} \quad (i \neq k).}
\label{eq:mass_mu_offdiagonal}
\end{equation}
Under $\mathrm{O}(K_q)$ this reduces to $-\beta_{ik}\Omega_{ik}\Lambda_{q_k} - \beta_{ki}\Lambda_{q_i}\Omega_{ki}^\top$.

\paragraph{Off-diagonal block caveats.}\label{sec:mass_block_caveats}
As the Hessian of the scalar functional $\mathcal{F}$, the block matrix $M_{\mu\mu}$ is automatically symmetric in the bilinear-form sense $M_{\mu\mu}^\top = M_{\mu\mu}$, equivalently $[M_{\mu\mu}]_{ik} = [M_{\mu\mu}]_{ki}^\top$, by equality of mixed partial derivatives (Schwarz's theorem). The block-transpose identity is exactly the condition for $M_{\mu\mu}$ to define a quadratic form $\tfrac12 \delta\mu^\top M_{\mu\mu} \delta\mu$ on the stacked variable $(\mu_i)_i$, the second-order term of $\mathcal{F}$; no stronger blockwise-equality condition is required, and an individual block $[M_{\mu\mu}]_{ik}$ need not be symmetric. Where the block-transpose identity appears not to hold in displayed expressions, the cause is that the displayed $[M_{\mu\mu}]_{ik}$ has been restricted to a subset of the contributions to the Hessian (for example a single sender-only or receiver-only contribution), and the symmetry is restored once both contributions are summed.

Row-asymmetric attention does not by itself destroy the scalar potential: fixed asymmetric weights still define the scalar $\sum_{ij}\beta_{ij}E_{ij}$, and state-dependent Gibbs weights in the canonical entropy-retaining objective reduce to $-\tau\sum_i\log Z_i$. In both cases the full receiver-plus-sender Hessian is symmetric, and the full flow is a gradient flow for every fixed $\beta$, reciprocal or not. A nonpotential running rule appears only after the additional receiver-only truncation of Eq.~\eqref{eq:receiver_only_field}, which discards the sender-row derivatives. For that optional field the bilinear form $M_{\mu\mu}$ is not the linearization of the dynamics; the linearization is the Jacobian of the truncated field, which is not symmetric and can have complex eigenvalues. The boxed forms of this section describe the full scalar $\mathcal{F}$ and do not describe the receiver-only field.

\paragraph{Covariance Sector}

For matrix-valued variables, we use the directional derivative convention
\begin{equation}
\frac{\partial^2 f}{\partial \Sigma \partial \Sigma}[V, W] = \lim_{\epsilon \to 0} \frac{1}{\epsilon}\left[\left.\frac{\partial f}{\partial \Sigma}\right|_{\Sigma + \epsilon W} - \left.\frac{\partial f}{\partial \Sigma}\right|_{\Sigma}\right][V],
\end{equation}
with the identity $\partial(\Sigma^{-1})/\partial \Sigma = -\Sigma^{-1} \otimes \Sigma^{-1}$.

\paragraph{Diagonal blocks $(i = k)$:} From the prior,
\begin{equation}
\frac{\partial^2 D_{\mathrm{KL}}(q_i \| p_i)}{\partial \Sigma_i \partial \Sigma_i}[V, W] = \frac{1}{2}\operatorname{tr}[\Lambda_{q_i} V \Lambda_{q_i} W].
\end{equation}
In tensor notation, $\partial^2 D_{\mathrm{KL}}(q_i \| p_i)/\partial \Sigma_i \partial \Sigma_i = \frac{1}{2}(\Lambda_{q_i} \otimes \Lambda_{q_i})$. From the interaction term as receiver and as sender, identical contributions arise \emph{at consensus} (see the remark below). The sensory term $\frac{1}{2}\operatorname{tr}(\Lambda_{o_i}\Sigma_i)$ is linear in $\Sigma_i$, so its second derivative vanishes. Therefore
\begin{equation}
\boxed{[M_{\Sigma\Sigma}]_{ii} = \frac{1}{2}(\Lambda_{q_i} \otimes \Lambda_{q_i}) \cdot \left(1 + \sum_k \beta_{ik} + \sum_j \beta_{ji}\right) \quad \text{at consensus.}}
\label{eq:mass_sigma_diagonal}
\end{equation}
The sensory precision $\Lambda_{o_i}$ contributes to the mean-sector stiffness but not to the covariance-sector stiffness.

\paragraph{Remark on the consensus simplification.} The factor $(1 + \sum_k \beta_{ik} + \sum_j \beta_{ji})$ in Eq.~\eqref{eq:mass_sigma_diagonal} arises because the sender contribution from $D_{\mathrm{KL}}(q_j \| \tilde{q}_i)$ involves both a negative term $-\frac{1}{2}\Lambda_{q_i} \otimes \Lambda_{q_i}$ from $\partial^2 \ln |\tilde{\Sigma}_i|/\partial \Sigma_i^2$ and a positive trace-term contribution that, away from consensus, equals $\frac{1}{2}\partial^2 \operatorname{tr}((\Omega_{ji}\Sigma_i\Omega_{ji}^\top)^{-1}\Sigma_j)/\partial\Sigma_i^2$ and couples $\Sigma_i$ to $\Sigma_j$. At consensus ($\Sigma_j = \Omega_{ji}\Sigma_i\Omega_{ji}^\top$) this trace-term contribution simplifies to $+\Lambda_{q_i} \otimes \Lambda_{q_i}$, and the sender net contribution collapses to $+\frac{1}{2}\Lambda_{q_i}\otimes\Lambda_{q_i}$, recovering Eq.~\eqref{eq:mass_sigma_diagonal}. Off consensus, the sender contribution carries an explicit $\Sigma_j$-dependent cross-coupling and the diagonal block does not factor into a multiple of $\Lambda_{q_i}\otimes\Lambda_{q_i}$; the boxed form above is the consensus-point restriction.

\paragraph{Off-diagonal blocks $(i \neq k)$:} From $D_{\mathrm{KL}}(q_i \| \tilde{q}_k)$, varying both $\Sigma_i$ and $\Sigma_k$, with $\tilde{\Lambda}_{q_k}$ defined by Eq.~\eqref{eq:precision_transport},
\begin{equation}
\boxed{\begin{aligned}
\relax [M_{\Sigma\Sigma}]_{ik} = &-\frac{1}{2}\beta_{ik}\left(\Omega_{ik}^{-\top}\Lambda_{q_k} \otimes \Omega_{ik}^{-\top}\Lambda_{q_k}\right) \\
&-\frac{1}{2}\beta_{ki}\left(\Omega_{ki}^{-\top}\Lambda_{q_i} \otimes \Omega_{ki}^{-\top}\Lambda_{q_i}\right)^\top \quad (i \neq k).
\end{aligned}}
\label{eq:mass_sigma_offdiagonal}
\end{equation}
The block is a pure Kronecker square of the two mean-sector transport factors. Each contribution is independent of the native covariance of the agent whose divergence produces it, a structural property any candidate form of this block must satisfy: agent $i$'s native covariance $\Sigma_i$ enters $D_{\mathrm{KL}}(q_i \| \tilde{q}_k)$ only linearly through the trace $\frac{1}{2}\operatorname{tr}(\tilde{\Lambda}_{q_k}\Sigma_i)$, so differentiating in $\Sigma_i$ removes $\Sigma_i$ entirely and leaves $\frac{1}{2}\tilde{\Lambda}_{q_k}$, and the subsequent differentiation in $\Sigma_k$ through the precision transport of Eq.~\eqref{eq:precision_transport} produces the Kronecker square $\Omega_{ik}^{-\top}\Lambda_{q_k} \otimes \Omega_{ik}^{-\top}\Lambda_{q_k}$, which depends on $\Sigma_k$ through $\Lambda_{q_k}$ but carries no residual $\Sigma_i$. By the same argument the sender contribution is independent of $\Sigma_k$ and depends on $\Sigma_i$ through $\Lambda_{q_i}$, so the assembled block depends on both precisions while each term is free of its own agent's covariance. The two-term structure mirrors the mean-sector off-diagonal block of Eq.~\eqref{eq:mass_mu_offdiagonal}: the first term is the receiver contribution from $D_{\mathrm{KL}}(q_i \| \tilde{q}_k)$ when $i$ attends to $k$, built from $A_{ik} = \Omega_{ik}^{-\top}\Lambda_{q_k}$; the second is the sender contribution from $D_{\mathrm{KL}}(q_k \| \tilde{q}_i)$ when $k$ also attends to $i$, built from $B_{ki} = \Omega_{ki}^{-\top}\Lambda_{q_i}$. The transpose on the second Kronecker factor is essential. The sender mixed partial is naturally assembled with the index ordering $(\Sigma_k, \Sigma_i)$, whereas the block $[M_{\Sigma\Sigma}]_{ik}$ requires the ordering $(\Sigma_i, \Sigma_k)$; since $\Omega_{ki}^{-\top}\Lambda_{q_i} \otimes \Omega_{ki}^{-\top}\Lambda_{q_i}$ is not symmetric, the index swap is exactly the vec-space transpose, and it is what makes the assembled block satisfy the Schwarz block-transpose identity $[M_{\Sigma\Sigma}]_{ik} = [M_{\Sigma\Sigma}]_{ki}^\top$. As in the mean sector, this identity holds for every fixed $\beta$, reciprocal or not, while an individual block $[M_{\Sigma\Sigma}]_{ik}$ need not be symmetric. Under $\mathrm{O}(K_q)$ the factors reduce to $\Omega_{ik}\Lambda_{q_k} \otimes \Omega_{ik}\Lambda_{q_k}$.

At the reciprocal-consensus point, the covariance-sector Hessian restricted to the symmetric tangent space on which $\Sigma$ lives is positive definite for all $\beta \geq 0$ and arbitrary precision heterogeneity. Take Regime I transports $\Omega_{ik} = \exp(\phi_i)\exp(-\phi_k)$, so that the cocycle holds and the transported neighbor $\tilde{q}_k$ coincides with $q_i$, with reciprocal $\beta$ and each prior equal to the native belief. Every relative entropy then sits at its global zero, contributing a positive-semidefinite Hessian, while the strictly convex prior term $\frac{1}{2}(\Lambda_{q_i} \otimes \Lambda_{q_i})$ on each diagonal block makes the total strictly positive definite. Four scope conditions are required. The result is established only at the consensus point; off consensus the $\Sigma_j$ cross-coupling re-enters and positive definiteness is not established there. It uses the frozen-$\beta$ (envelope-theorem) convention of the differentiation paragraph above. It requires a global consensus through the cocycle, not merely pairwise reciprocity; the full $K_q^2 \times K_q^2$ vec operator can be indefinite on the non-symmetric complement, which is not the physical object, since $\Sigma$ is constrained to the symmetric subspace. And the point is a critical point of $\mathcal{F}$ only when the sensory term is absent: the sensory term adds no covariance curvature but contributes the constant covariance gradient $\tfrac12\Lambda_{o_i}$, so with sensory evidence the statement concerns the flat-chart second variation at a noncritical point, with the chart dependence described at the start of this subsection.

\paragraph{Cross Mean-Covariance Blocks}

The prior term gives $\partial^2 D_{\mathrm{KL}}(q_i \| p_i)/\partial \mu_i \partial \Sigma_i = 0$. The sensory free energy decomposes into a term quadratic in $\mu_i$ and a term linear in $\Sigma_i$; these are independent, so $[C^{\mu\Sigma}]_{ii}^{\mathrm{sensory}} = 0$.

For the interaction term, from $\partial D_{\mathrm{KL}}(q_i \| \tilde{q}_k)/\partial \mu_i = \tilde{\Lambda}_{q_k}(\mu_i - \tilde{\mu}_k)$ with $\tilde{\Lambda}_{q_k} = \Omega_{ik}^{-\top}\Lambda_{q_k}\Omega_{ik}^{-1}$ (Eq.~\eqref{eq:precision_transport}), varying $\Sigma_k$ gives $\delta\tilde{\Lambda}_{q_k} = -\Omega_{ik}^{-\top}\Lambda_{q_k} V \Lambda_{q_k}\Omega_{ik}^{-1}$ for $V = \delta\Sigma_k$, and therefore
\begin{equation}
\frac{\partial^2 D_{\mathrm{KL}}(q_i \| \tilde{q}_k)}{\partial \mu_i \partial \Sigma_k}[V] = -\Omega_{ik}^{-\top}\Lambda_{q_k} V \Lambda_{q_k} \Omega_{ik}^{-1} (\mu_i - \tilde{\mu}_k).
\end{equation}
In components,
\begin{equation}
\frac{\partial^2 D_{\mathrm{KL}}}{\partial(\mu_i)_a \partial(\Sigma_k)_{bc}} = -[\Omega_{ik}^{-\top}\Lambda_{q_k}]_{ab}[\Lambda_{q_k}\Omega_{ik}^{-1}(\mu_i - \tilde{\mu}_k)]_c.
\end{equation}
Under $\mathrm{O}(K_q)$ this reduces to $-\Omega_{ik}\Lambda_{q_k} V \Lambda_{q_k}\Omega_{ik}^\top(\mu_i - \tilde{\mu}_k)$. The cross block vanishes at consensus ($\mu_i = \tilde{\mu}_k$):
\begin{equation}
\boxed{[C^{\mu\Sigma}]_{ik} = 0 \quad \text{when } \mu_i = \tilde{\mu}_k.}
\label{eq:cross_block}
\end{equation}

\subsection{Precision as Stiffness and the Fisher Metric as Friction}
\label{sec:stiffness_friction}

The mean-sector stiffness of agent $i$, in the frozen-$\beta$ convention, is
\begin{equation}
M_i = \underbrace{\bar{\Lambda}_{p_i}}_{\text{prior}} + \underbrace{\sum_k \beta_{ik}\tilde{\Lambda}_{q_k}}_{\text{receiver}} + \underbrace{\sum_j \beta_{ji}\Lambda_{q_i}}_{\text{sender}} + \underbrace{\Lambda_{o_i}}_{\text{sensory}}.
\label{eq:effective_mass}
\end{equation}
The prior block resists deviation from the prior $p_i$, the receiver block resists deviation from the transported neighbors that agent $i$ attends to, the sender block resists moves that would disturb the agents attending to $i$, and the sensory block anchors the agent to its observations. The stiffness is dimensionless on the representational space; no identification with a physical mass is made.

\paragraph{Constant versus state-dependent self-coupling.} The boxed stiffness blocks of this section take the prior self-coupling weight at its canonical default $\alpha_i = 1$; the prior term of Eq.~\eqref{eq:extended_free_energy} carries unit weight. Under the state-dependent generalization of Section~\ref{sec:state_dependent_precision}, where $\alpha_i^*(c) = c_0/(b_0 + D_{\mathrm{KL}}(q_i \| p_i))$ is recomputed from the current belief at each inner step, the prior block scales as $\alpha_i^* \bar{\Lambda}_{p_i}$. Because $\alpha_i^*$ co-varies with the belief, the second variation acquires a rank-one residual that the first-order envelope cancellation of Eq.~\eqref{eq:alpha_envelope_cancel_itfb} does not remove,
\begin{equation}
[M_{\mu\mu}]_{ii}^{\text{prior}} = \alpha_i^* \bar{\Lambda}_{p_i} - \frac{(\alpha_i^*)^2}{c_0}\left(\bar{\Lambda}_{p_i}\Delta\mu_i\right)\left(\bar{\Lambda}_{p_i}\Delta\mu_i\right)^\top, \qquad \Delta\mu_i = \mu_i - \bar{\mu}_i,
\end{equation}
which vanishes only at $\mu_i = \bar{\mu}_i$. This is the prior-sector analog of the off-consensus rank-one stiffness retained for the interaction term in the first-variation analysis above. For a mean displacement at matched covariance, where $D_{\mathrm{KL}}(q_i \| p_i) = \tfrac{1}{2}\Delta\mu_i^\top \bar{\Lambda}_{p_i}\Delta\mu_i$, the regularized prior block stays positive definite until $D_{\mathrm{KL}} = b_0$, whereas the unregularized weight $\alpha_i^* D_{\mathrm{KL}}$ alone would lose definiteness already at $D_{\mathrm{KL}} = b_0/3$; the log-barrier delays it. Beyond that bound the prior contribution to the stiffness is indefinite. The boxed forms of this section are therefore the quasi-static $\alpha_i = 1$ specialization.

\paragraph{Friction and relaxation.} The flow of Eq.~\eqref{eq:gauge_natural_gradient} is first order. In the Gaussian sector it reads
\begin{equation}
\eta_b^{-1}G_b\dot{\xi}_b = -\nabla_{\xi_b}\mathcal{F},
\qquad
G_\mu = \Lambda_q,
\qquad
G_\Sigma[V, W] = \tfrac12\operatorname{tr}(\Sigma^{-1}V\Sigma^{-1}W),
\label{eq:friction_form}
\end{equation}
for each block $b \in \{\mu, \Sigma\}$ with its rate $\eta_b$, where $G_b$ is the Fisher metric of the block. This is the form $\gamma\dot{q} = -\nabla V$ of an overdamped system with friction tensor $\gamma = \eta^{-1}G$ and mobility $\eta G^{-1}$, and it is the Rayleigh form $\partial\mathcal{R}/\partial\dot{\xi} = -\nabla\mathcal{F}$ with dissipation function
\begin{equation}
\mathcal{R}(\dot{\xi}) = \sum_b \frac{1}{2\eta_b}G_b(\dot{\xi}_b, \dot{\xi}_b),
\qquad
\frac{d\mathcal{F}}{dt} = \nabla\mathcal{F}\cdot\dot{\xi} = -2\mathcal{R}.
\label{eq:rayleigh_dissipation}
\end{equation}
With a common rate $\eta_b = \eta$ the squared Fisher speed is $G(\dot{\xi}, \dot{\xi}) = (d\vartheta/dt)^2 = 2\eta\mathcal{R}$, which is the clock-dissipation identity of Eq.~\eqref{eq:clock_dissipation} for an autonomous single-rate flow. The block-descent law and the Fisher clock are established in general in the companion theory~\cite[Chs.~7 and 8]{DennisGaugeVFE2026}. Linearized about a critical point $\xi^*$, the flow is $\delta\dot{\xi} = -\eta G^{-1}M\delta\xi$. Since $G$ is positive definite and $M$ is symmetric, $G^{-1}M$ is similar to $G^{-1/2}MG^{-1/2}$, so its spectrum is real and the linearized flow has no oscillatory modes; complex eigenvalues arise only in the receiver-only truncation of Eq.~\eqref{eq:receiver_only_field}. For an isolated agent at $q_i = p_i$ the stiffness and the friction coincide, $[M_{\mu\mu}]_{ii} = \bar{\Lambda}_{p_i} = \Lambda_{q_i} = G_\mu$ and $[M_{\Sigma\Sigma}]_{ii} = \tfrac12\Lambda_{q_i}\otimes\Lambda_{q_i} = G_\Sigma$, so the linearized mean and covariance relaxation rates are exactly $\eta_\mu$ and $\eta_\Sigma$, independent of the precision. Under a Euclidean gradient flow $\dot{\xi} = -\eta\nabla\mathcal{F}$ the rates would instead be $\eta$ times the eigenvalues of $M$, so a precise agent would relax faster. In neither flow is a high-precision agent slow; large precision produces large stiffness and small static compliance $M^{-1}$.

\paragraph{High-precision agents.} Consider an agent with very sharp prior self-localization, $\bar{\Sigma}_i = \delta I$ with $\delta \ll 1$ a dimensionless precision scale to which no physical units are assigned (the macroscopic-object example of Appendix~\ref{sec:rock_derivations}). In the isolated-agent limit of Eq.~\eqref{eq:effective_mass}, where the interaction blocks switch off ($\beta_{ik} = \beta_{ji} = 0$) and the sensory block is negligible ($\Lambda_{o_i} \to 0$), $M_i$ reduces to the prior precision $\delta^{-1}I$; the interaction and sensory terms add further stiffness when the agent is embedded. Such an agent has large stiffness and compliance of order $\delta$, so a bounded force displaces it by an amount of order $\delta$; under the natural-gradient flow its linear relaxation rate is still $\eta$.

The stiffness above runs entirely on the state-fiber prior precision $\bar{\Lambda}_{p_i}$ and is therefore specific to the prior-pullback tier in the sense of Section~\ref{sec:three_tiers}. Under the structural-tier reading of Section~\ref{sec:gravity_pullback}, perceived geometry is carried instead by the model-fiber generative-model precision $\Sigma_{s,i}^{-1}$, so the stiffness of this section and the pulled-back geometry are distinct objects at that tier; relating them would require a dynamical link between $\bar{\Lambda}_{p_i}$ and $\Sigma_{s,i}^{-1}$ that the present construction does not supply.

\subsection{A Declared Kinetic-Metric Postulate}
\label{sec:velocity_quadratic}

\paragraph{Postulate (kinetic metric).} In addition to $\mathcal{F}$ and its first-order flow, one may postulate a kinetic energy $\mathcal{T} = \tfrac12 \dot{\xi}^\top M_K \dot{\xi}$ with inertia tensor $M_K$ and adopt the second-order equation of motion
\begin{equation}
M_K\ddot{\xi} + \eta^{-1}G\dot{\xi} = -\nabla\mathcal{F},
\label{eq:second_order_postulate}
\end{equation}
the Euler-Lagrange equation of $\mathcal{T} - \mathcal{F}$ with the Rayleigh dissipation of Eq.~\eqref{eq:rayleigh_dissipation}. Its overdamped limit, in which the inertial times are short compared with the relaxation times, is the natural-gradient flow of Eq.~\eqref{eq:friction_form}. Neither $\mathcal{F}$ nor the flow of Eq.~\eqref{eq:gauge_natural_gradient} supplies $M_K$. The inertial reading of precision that the postulate expresses has status S (Table~\ref{tab:claim_status}), and no other result of the paper uses Eq.~\eqref{eq:second_order_postulate}. The natural candidate takes $M_K = M$, the frozen-$\beta$ stiffness at consensus. Its consequences are fixed by linearizing about a critical point. Without dissipation the linearized equation is $M\delta\ddot{\xi} = -M\delta\xi$, so $\delta\ddot{\xi} = -\delta\xi$: every linear mode has the same frequency, $\omega = 1$ in flow-time units, whatever the precision. With the dissipation retained, an isolated agent at $q_i = p_i$, where $M = G$, obeys $\delta\ddot{\xi} + \eta^{-1}\delta\dot{\xi} + \delta\xi = 0$, with damping ratio $1/(2\eta)$; the motion is overdamped for $\eta < 1/2$, and for $\eta \ll 1/2$ the slow root is $-\eta$, the natural-gradient rate. Under the postulate $M_K = M$ the stiffness and the inertia coincide, so no precision-dependent frequency, and in particular no scaling $\omega^2 \propto 1/\bar{M}_i$, follows from it. A precision-dependent frequency requires an inertia tensor that is not the stiffness tensor, which the construction does not supply. For example, the different postulate $M_K = \mathrm{diag}_i(\bar{\Lambda}_{p_i})$ in the mean sector, with $K_c := M_{\mu\mu} - M_K$ the interaction-plus-sensory stiffness, gives $\omega^2 = \mathrm{eig}(M_K^{-1}M_{\mu\mu}) = 1 + \mathrm{eig}(M_K^{-1}K_c)$, a testable relation between frequencies and coupling-plus-sensory stiffness, not an identity; we do not adopt it. A dynamical test of any such postulate requires a trajectory generated by Eq.~\eqref{eq:second_order_postulate}, for example underdamped ringing, which the first-order dynamics used throughout this paper cannot produce, because its linearized spectrum is real; a static-response measurement of the compliance $M^{-1}$ tests the stiffness and says nothing about $M_K$.

Under the postulate $M_K = M$ the kinetic energy takes the form
\begin{equation}
\mathcal{T} = \frac{1}{2}\dot{\mu}^\top M_{\mu\mu} \dot{\mu} + \frac{1}{2}M_{\Sigma\Sigma}[\dot{\Sigma}, \dot{\Sigma}] + \frac{1}{2} G_\phi(\phi)[\dot{\phi}, \dot{\phi}],
\label{eq:full_kinetic}
\end{equation}
where the frame-sector coefficient $G_\phi$ is a further choice, fixed by cases. For $G=\mathrm{SO}(K_q)$ with $K_q\ge2$, the frame $\phi\in\mathfrak{so}(K_q)$ carries the bi-invariant inner product $G_\phi[\dot{\phi},\dot{\phi}]=\langle\dot{\phi},\dot{\phi}\rangle_{\mathfrak{g}}=-\operatorname{tr}(\dot{\phi}^2)$. When $K_q\ge3$, $\mathfrak{so}(K_q)$ is compact semisimple and this inner product is $1/(K_q-2)$ of the negative Killing form, since $B(X,X)=(K_q-2)\operatorname{tr}(X^2)$. At $K_q=2$, $\mathfrak{so}(2)$ is abelian and its Killing form vanishes, so the positive Ad-invariant negative trace form is a separate choice, not a Killing-form construction. Central or abelian factors in any other compact algebra require the same separate choice of positive Ad-invariant inner product. For the general $\mathrm{GL}(K_q)$ case the Killing form on $\mathfrak{gl}(K_q)$ is degenerate (its radical is the center) and indefinite on the semisimple part $\mathfrak{sl}(K_q)$, and $G_\phi$ would be supplied instead by the position-dependent pullback $\mathcal{G}_{\mathrm{pull}}$ of the ambient Frobenius metric through the exponential map, Eq.~\eqref{eq:pullback_metric}, under the $\mathrm{dexp}^{L}_\phi$-regularity hypothesis recorded at Eq.~\eqref{eq:gauge_natural_gradient_def}. That form is flat, being the pullback of a flat metric by a local diffeomorphism, and its kinetic term $\tfrac12\mathcal{G}_{\mathrm{pull}}(\dot{\phi}, \dot{\phi}) = \tfrac12\|\dot{U}\|_F^2$ is neither left nor right invariant. Under the common right translation $U_i \to U_i g$, which leaves every Regime I transport and hence $\mathcal{F}$ unchanged, the ratio $\|\dot{U}g\|_F^2/\|\dot{U}\|_F^2 = \operatorname{tr}(\dot{U}gg^\top\dot{U}^\top)/\operatorname{tr}(\dot{U}\dot{U}^\top)$ ranges over the interval between the smallest and largest eigenvalues of $gg^\top$, so the term is invariant for every velocity only when $g$ is orthogonal. Since $\mathrm{GL}^+(K_q)$ admits no bi-invariant Riemannian metric~\cite{Milnor1976}, any frame-sector kinetic term for noncompact frames breaks part of the redundancy that $\mathcal{F}$ respects, and in Regime I, where the frames are removable passive re-expressions (Section~\ref{sec:one_rung}), the frame sector of Eq.~\eqref{eq:full_kinetic} is an additional, frame-dependent postulate.

The second term of Eq.~\eqref{eq:full_kinetic} couples covariance evolution to the covariance-sector stiffness, which at an isolated equilibrium is the Fisher (affine-invariant) metric on $\mathbb{S}^+_{K_q}$. The inter-agent off-diagonal blocks contribute
\begin{equation}
\mathcal{T}_{\mathrm{couple}} = \sum_{i < k}\left[-\beta_{ik}\dot{\mu}_i^\top \Omega_{ik}^{-\top}\Lambda_{q_k}\dot{\mu}_k - \beta_{ki}\dot{\mu}_i^\top \Lambda_{q_i}\Omega_{ki}^{-1} \dot{\mu}_k\right],
\end{equation}
which under the postulate couples the accelerations of the agents. By the block-transpose symmetry of the off-diagonal block caveats paragraph above, this quadratic form is well defined for any $\beta$, reciprocal or not. Equality of the two displayed coefficient blocks is a separate requirement: it holds if and only if $\beta_{ik} = \beta_{ki}$ and the transport identity $\Omega_{ik}^{-\top}\Lambda_{q_k} = \Lambda_{q_i}\Omega_{ki}^{-1}$ hold, and under orientation reversal $\Omega_{ki} = \Omega_{ik}^{-1}$ the latter reduces to the consensus condition $\Omega_{ik}^{-\top}\Lambda_{q_k}\Omega_{ik}^{-1} = \Lambda_{q_i}$ (the transported precisions agree), which reciprocity alone does not enforce.
